% GENERATED FILE. Edit canonical lessons, then run npm run generate:books.
% canonical-source-hash: fnv1a64:d869c5a5d72dcc90
% canonical-lessons: IT-C366-poverta, IT-C366-disoccupazione, IT-C366-discriminazione, IT-C366-discriminare, IT-C366-immigrazione

\chapter{Problemi sociali}
\label{ch:it-problemi-sociali}
\hlchaptermodality{366}

\begin{chapteropening}
\textbf{By the end of this chapter:} \emph{I can talk about social problems: poverty, unemployment, discrimination and immigration.}

Complete the last lesson of chapter 366: I can talk about social problems: poverty, unemployment, discrimination and immigration.
\end{chapteropening}

\section[la povertà]{la povertà — the poverty}

\label{lesson:IT-C366-poverta}



\begin{quote}
Before the new one: say the Italian for the prejudice, then the Italian for the stereotype.
\end{quote}

\subsection*{You'll want to know: la povertà}
\textbf{la povertà} — \textquotedblleft{}the poverty\textquotedblright{}.

\begin{quote}
Molte famiglie vivono in povertà.
\end{quote}

\begin{quote}
\emph{Many families live in poverty.}
\end{quote}

Built on \emph{povero}.

\begin{cousinweb}
\emph{povertà} is inherited from Latin \emph{paupertās, paupertātis} \textquotedblleft{}poverty\textquotedblright{} (Lewis \& Short), from \emph{pauper} \textquotedblleft{}poor\textquotedblright{}.
\end{cousinweb}

\begin{grammarlens}[title={What you've built}]
One more word for telling, explaining and describing at B1.
\end{grammarlens}

\subsection*{Your turn}
\begin{itemize}
\raggedright
  \item \emph{Say it:} \emph{la povertà}
  \item \emph{Say it:} \emph{la povertà}, once more
  \item \emph{Say it:} a sentence of your own with \emph{la povertà}
  \item \emph{Recall:} say the Italian for the prejudice, then the Italian for the stereotype, then say \emph{la povertà} again
\end{itemize}

\begin{tcolorbox}[breakable,colback=teal!4,colframe=teal!35!black,title={Before you move on}]
What does la povertà mean? (\textquotedblleft{}The poverty\textquotedblright{}.) Say it once more.
\end{tcolorbox}
\section[la disoccupazione]{la disoccupazione — the unemployment}

\label{lesson:IT-C366-disoccupazione}



\begin{quote}
Before the new one: say the Italian for the stereotype, then the Italian for the poverty.
\end{quote}

\subsection*{You'll want to know: la disoccupazione}
\textbf{la disoccupazione} — \textquotedblleft{}the unemployment\textquotedblright{}.

\begin{quote}
La disoccupazione tra i giovani è alta.
\end{quote}

\begin{quote}
\emph{Unemployment among young people is high.}
\end{quote}

Built on \emph{disoccupato}.

\begin{cousinweb}
\emph{disoccupazione} is built on \emph{disoccupato} \textquotedblleft{}unemployed\textquotedblright{}, from \emph{dis-} and \emph{occupato}, from Latin \emph{occupāre} \textquotedblleft{}to seize, to fill\textquotedblright{} (Lewis \& Short), built on \emph{capere} \textquotedblleft{}to take\textquotedblright{}.
\end{cousinweb}

\begin{grammarlens}[title={What you've built}]
One more word for telling, explaining and describing at B1.
\end{grammarlens}

\subsection*{Your turn}
\begin{itemize}
\raggedright
  \item \emph{Say it:} \emph{la disoccupazione}
  \item \emph{Say it:} \emph{la disoccupazione}, once more
  \item \emph{Say it:} a sentence of your own with \emph{la disoccupazione}
  \item \emph{Recall:} say the Italian for the stereotype, then the Italian for the poverty, then say \emph{la disoccupazione} again
\end{itemize}

\begin{tcolorbox}[breakable,colback=teal!4,colframe=teal!35!black,title={Before you move on}]
What does la disoccupazione mean? (\textquotedblleft{}The unemployment\textquotedblright{}.) Say it once more.
\end{tcolorbox}
\section[la discriminazione]{la discriminazione — the discrimination}

\label{lesson:IT-C366-discriminazione}



\begin{quote}
Before the new one: say the Italian for the poverty, then the Italian for the unemployment.
\end{quote}

\subsection*{You'll want to know: la discriminazione}
\textbf{la discriminazione} — \textquotedblleft{}the discrimination\textquotedblright{}.

\begin{quote}
La discriminazione sul lavoro è contro la legge.
\end{quote}

\begin{quote}
\emph{Discrimination at work is against the law.}
\end{quote}

\emph{Discriminazione contro le donne}.

\begin{cousinweb}
\emph{discriminazione} is a learned borrowing from Late Latin \emph{discrīminātiō}, a term of rhetoric for setting opposite ideas side by side (Lewis \& Short), from \emph{discrīmen} \textquotedblleft{}a dividing line\textquotedblright{}, built on \emph{discernere} and \emph{cernere} \textquotedblleft{}to sift, to separate\textquotedblright{}.
\end{cousinweb}

\begin{grammarlens}[title={What you've built}]
One more word for telling, explaining and describing at B1.
\end{grammarlens}

\subsection*{Your turn}
\begin{itemize}
\raggedright
  \item \emph{Say it:} \emph{la discriminazione}
  \item \emph{Say it:} \emph{la discriminazione}, once more
  \item \emph{Say it:} a sentence of your own with \emph{la discriminazione}
  \item \emph{Recall:} say the Italian for the poverty, then the Italian for the unemployment, then say \emph{la discriminazione} again
\end{itemize}

\begin{tcolorbox}[breakable,colback=teal!4,colframe=teal!35!black,title={Before you move on}]
What does la discriminazione mean? (\textquotedblleft{}The discrimination\textquotedblright{}.) Say it once more.
\end{tcolorbox}
\section[discriminare]{discriminare — to discriminate (against)}

\label{lesson:IT-C366-discriminare}



\begin{quote}
Before the new one: say the Italian for the unemployment, then the Italian for the discrimination.
\end{quote}

\subsection*{You'll want to know: discriminare}
\textbf{discriminare} — \textquotedblleft{}to discriminate (against)\textquotedblright{}.

\begin{quote}
Non si può discriminare una persona per la sua religione.
\end{quote}

\begin{quote}
\emph{You cannot discriminate against a person because of their religion.}
\end{quote}

Stress: \emph{discrìmino}. The old sense is to tell apart.

\begin{cousinweb}
\emph{discriminare} is a learned borrowing from Latin \emph{discrīmināre} \textquotedblleft{}to divide, to separate\textquotedblright{} (Lewis \& Short), from \emph{discrīmen} \textquotedblleft{}a dividing line\textquotedblright{}, built on \emph{cernere} \textquotedblleft{}to sift\textquotedblright{}.
\end{cousinweb}

\begin{grammarlens}[title={What you've built}]
One more word for telling, explaining and describing at B1.
\end{grammarlens}

\subsection*{Your turn}
\begin{itemize}
\raggedright
  \item \emph{Say it:} \emph{discriminare}
  \item \emph{Say it:} \emph{discriminare}, once more
  \item \emph{Say it:} a sentence of your own with \emph{discriminare}
  \item \emph{Recall:} say the Italian for the unemployment, then the Italian for the discrimination, then say \emph{discriminare} again
\end{itemize}

\begin{tcolorbox}[breakable,colback=teal!4,colframe=teal!35!black,title={Before you move on}]
What does discriminare mean? (\textquotedblleft{}To discriminate (against)\textquotedblright{}.) Say it once more.
\end{tcolorbox}
\section[l'immigrazione]{l'immigrazione — the immigration}

\label{lesson:IT-C366-immigrazione}



\begin{quote}
Before the new one: say the Italian for the discrimination, then the Italian for to discriminate.
\end{quote}

\subsection*{You'll want to know: l'immigrazione}
\textbf{l'immigrazione} — \textquotedblleft{}the immigration\textquotedblright{}.

\begin{quote}
Il governo vuole cambiare le leggi sull'immigrazione.
\end{quote}

\begin{quote}
\emph{The government wants to change the immigration laws.}
\end{quote}

\emph{Un immigrato} comes in; \emph{un emigrato} has left.

\begin{cousinweb}
\emph{immigrazione} is built on \emph{immigrare} \textquotedblleft{}to immigrate\textquotedblright{}, a learned borrowing from Latin \emph{immigrāre} \textquotedblleft{}to move into\textquotedblright{} (Lewis \& Short), built from \emph{in} and \emph{migrāre} \textquotedblleft{}to move house\textquotedblright{}.
\end{cousinweb}

\begin{grammarlens}[title={What you've built}]
One more word for telling, explaining and describing at B1.
\end{grammarlens}

\subsection*{Your turn}
\begin{itemize}
\raggedright
  \item \emph{Say it:} \emph{l'immigrazione}
  \item \emph{Say it:} \emph{l'immigrazione}, once more
  \item \emph{Say it:} a sentence of your own with \emph{l'immigrazione}
  \item \emph{Recall:} say the Italian for the discrimination, then the Italian for to discriminate, then say \emph{l'immigrazione} again
\end{itemize}

\begin{tcolorbox}[breakable,colback=teal!4,colframe=teal!35!black,title={Before you move on}]
What does l'immigrazione mean? (\textquotedblleft{}The immigration\textquotedblright{}.) Say it once more.
\end{tcolorbox}
